\section{Why a common-exponent counting argument is insufficient}
\label{sec:critical-counting-obstruction}

Let $m=2^d$ and
\[
 E_{N,L}=C\left[\frac{N}{L-1}+(L-1)\log(2N)\right],
\]
with $C$ as in the finite-width universal upper bound. The target
events $A_{N,f}$ are disjoint and their union is the global-margin
event $A_N$. Therefore
\[
 \sum_f\mu_{N,1}(A_{N,f})
 \le 2\exp(-NV_*+E_{N,L}).
\]
This gives an immediate obstruction to transferring the common
variational rate into uniform finite-width target masses.

\begin{proposition}[Unavoidable function-specific logarithmic overhead]
If a number $R$ satisfies
\[
 \mu_{N,1}(A_{N,f})\ge\exp(-NV_*-R)
 \qquad\text{for every Boolean }f,
\]
then necessarily
\[
 R\ge m\log2-\log2-E_{N,L}.
\]
In particular, in a regime where $E_{N,L}=o(m)$, a uniform lower
bound with $R=o(m)$ is impossible.
\end{proposition}

\begin{proof}
Sum the asserted lower bound over the $2^m$ sign patterns and compare
with the displayed upper bound on their union.
\end{proof}

More generally, for every $r\in\mathbb R$, the number of functions
whose target event has probability at least $e^{-NV_*-r}$ is at most
$2 e^{E_{N,L}+r}$. Thus the common limiting exponent must conceal
substantial function-specific finite-width terms. Their presence is
not an artifact of a particular coupling or concentration proof.

This observation does not establish good learning: the missing terms
could favor simple targets, or could resemble a nearly uniform prior
over truth tables. It also does not establish bad learning. A posterior
sample lower bound requires a target-relative comparison that retains
these terms; a shared leading exponent alone cannot supply one.
